% 1. Introduction: A Bottle in the Roof
\section{Introduction: A Bottle in the Roof}
\label{sec:introduction}

In 2002, during one of the electricity blackouts that kept hitting his part of
Brazil, a mechanic in Uberaba named Alfredo Moser had an idea that needed
almost nothing. He took a clear two-litre plastic bottle, filled it with water,
added two capfuls of bleach so the water would not turn green, and pushed it
through a hole in a tin roof, sealing the gap around it. In daylight the water
bends the sunlight that falls on the top of the bottle and spreads it into the
room below. A dark shack gets a lamp that has no wire, no switch and no bill.
Moser liked to say that an engineer had measured it at ``more or less
\mooserWattsLow{} to \mooserWattsHigh{} watts''
\parencite{zobel2013bbc}. That is his estimate, not an independent test, but
anyone who has seen one lit will understand why he was proud of it.

The idea travelled. In June 2011 the MyShelter Foundation in the Philippines
began making and installing the lamps under the name Liter of Light, and by
the summer of 2013 the foundation reported them in \literHomesPhilippines{}
Philippine homes and in about \literCountries{} other countries
\parencite{zobel2013bbc}. A device that costs almost nothing had crossed an
ocean.

\begin{pullquote}
The light bulb needed a power grid. The bottle light needed a roof and the sun.
\end{pullquote}

I keep coming back to that contrast because it captures something I think is
true and under-appreciated. The light bulb is a wonderful invention, but it is
the visible tip of an enormous system: power stations, transmission lines,
meters, utilities, a billing relationship and a building with wiring. Where
that system exists, the light bulb is the obvious answer to a dark room. Where
it does not, the obvious answer is not available, and someone has to ask the
question again from the beginning. Moser did not ask ``how do I get
electricity?'' He asked ``how do I get light into this room during the day?'',
and that question has a much cheaper answer.

This paper is about that move: what happens when people lack the resources an
existing solution takes for granted, and are forced to rethink the problem.
My argument is modest. Scarcity does not produce \emph{more} innovation.
Most of the time it produces less, or none, and simply means going without.
But sometimes it produces a \emph{different kind} of innovation: solutions
that are cheaper, simpler and more widely usable than the original, because
they were never allowed to assume the original's surroundings. A few of these
solutions then travel back to the places that had the resources all along.

\begin{keyidea}
Scarcity is not a source of more innovation. It is, occasionally, a source of
a different kind: solutions designed without the surrounding system that the
standard answer quietly depends on.
\end{keyidea}

The interesting question is therefore not \emph{whether} scarcity breeds
invention, which it plainly often does not, but \emph{when}. Under what
conditions does a missing grid lead to a bottle light or a solar kit, rather
than to another decade of kerosene? That is the question I try to answer, and
I should say plainly at the start that I can only answer it partly.

I want to be honest about the problems with this argument before making it,
because they are serious and I return to each of them in
\cref{sec:against}. The first is survivorship bias: we tell the story of the
bottle light because it worked, and we do not tell the stories of the millions
of homes where scarcity just meant darkness. The second is that much of what
looks like frugal invention runs on hidden abundance. A cheap solar lantern is
cheap because of semiconductor fabs, battery factories and global shipping,
none of which were built under scarcity. The third is that the fixes have real
limits. The bottle light does nothing at night. Today about
\sdgWithoutElecWorld{} people still have no electricity at all, \sdgWithoutElecSSA{}
of them in sub-Saharan Africa \parencite{esmap2026sdg7}, and no amount of
ingenuity with bottles changes that. The fourth is the risk of romanticising
hardship: praising the resourcefulness of poor communities can slide into
excusing the underinvestment that made resourcefulness necessary. I would like
this paper to do the opposite.

The rest of the paper goes as follows. \Cref{sec:terms} defines terms and sets
out the mechanism in three parts. \Cref{sec:case-for} makes the case for it
using five bodies of evidence: off-grid energy, mobile money, low-cost health
care, the psychology of constraints and a few historical echoes. Each part
comes with its own counterpoint. \Cref{sec:against} argues against the whole
thing as hard as I can. \Cref{sec:synthesis} tries to put the pieces back
together, and \cref{sec:conclusion} restates the thesis more carefully than I
have stated it here.
